\section*{Chapter \ref{ch:ll:the-order-book-builder} --- The Order Book Builder}

\begin{solution}{exo:ll:the-order-book-builder:1}
The base is $50.00 - 2\,048 \times 0.01 = 29.52$; 49.73 sits at index $(49.73 - 29.52)/0.01 = 2\,021$. The middle half is
indices 1\,024 to 3\,071, prices 39.76 to 60.23: the ladder stays put while both best prices remain in that range.
\end{solution}

\begin{solution}{exo:ll:the-order-book-builder:2}
$65\,536 \times 16 \times 2 = 2$~MiB, exactly the size of the laptop's second-level cache (chapter~3): it fits only if nothing
else does. In practice the touch and the few levels near it are what stays in cache.
\end{solution}

\begin{solution}{exo:ll:the-order-book-builder:3}
The replace removes the order of 500 at 20.10 (its level loses 500) and adds a new order of 300 at 20.11 under the new
reference, at the back of that level's queue: it has lost its priority, as the specification's new reference says.
\end{solution}

\begin{solution}{exo:ll:the-order-book-builder:4}
Before: slot 5 holds 1\,024, slot 6 2\,048, slot 7 3\,072. Erasing 1\,024 empties slot 5; 2\,048 (home 5) moves back to 5 and
3\,072 (home 5) to 6; slot 7 is emptied. A tombstone would have left slot 5 marked deleted: every lookup of 2\,048 and 3\,072
would probe through it until the table is rebuilt, and after a day of cancels, probe runs would grow long.
\end{solution}

\begin{solution}{exo:ll:the-order-book-builder:5}
$\sigma = 40 \times 0.02 / 0.01 = 80$ ticks a day. With 1\,024 ticks, $a = 256$: $(80/256)^2 \approx 0.098$, about one
recentring in ten days; with 4\,096, $a = 1\,024$: $0.0061$, one in 160 days.
\end{solution}

\begin{solution}{exo:ll:the-order-book-builder:6}
It removes the order (reducing its level by the 200 it holds) and counts the event as inconsistent: a reduction larger than
the order means an event was lost or duplicated upstream, and the monitoring should raise it with the \omterm{def:ll:the-feed-handler:handler}{feed handler}'s gap
counters.
\end{solution}

\begin{solution}{exo:ll:the-order-book-builder:7}
Keep, per side, an array of 64-bit words with one bit per ladder slot, set when a level becomes non-empty and cleared when it
empties. The next best bid below index $i$ is found word by word with count-leading-zeros on the masked word, the next best ask
with count-trailing-zeros: a few instructions per 64 empty slots instead of 64 loads. Recentring rebuilds the bitmaps with the
levels.
\end{solution}

\begin{solution}{exo:ll:the-order-book-builder:8}
A quiet day says nothing about the day the price moves: a 512-tick ladder ($a = 128$) recentres once per 128 ticks of travel,
and on a small-tick stock a thirty percent day is hundreds of thousands of ticks, thousands of recentrings (the measurement's
256-tick ladder needed almost twenty thousand). Size the ladder for the largest move expected, and test it on that move.
\end{solution}

\begin{solution}{pb:ll:the-order-book-builder:1}
\begin{enumerate}
\item $100 \times 0.015 / 10^{-4} = 15\,000$ ticks.
\item $(15\,000/1\,024)^2 \approx 215$ a day with 4\,096 ticks; $(15\,000/16\,384)^2 \approx 0.84$ with 65\,536.
\item One a day when $a = \sigma$: $4 \times 15\,000 = 60\,000$ ticks.
\item The ladder's cost comes from the days the price travels far; an ordinary day's rate hides them.
\item 300\,000 ticks; a steady trend causes $300\,000/1\,024 \approx 293$.
\item 256: each recentring centres the ladder on the mid, which has already passed the edge, so the next one needs more than
$a$ ticks of travel.
\item The book itself is wider than the middle half of a 256-tick ladder (orders up to 300 ticks from the mid), so the best
prices leave it at every jump of the touch, not only when the mid travels: 19\,744 recentrings.
\item $(12.8 - 4.4)~\text{ms}/(256 - 15) \approx \qty{35}{\micro\second}$ each: a pass over the order pool and the ladder.
\item 65\,536 ticks: fifteen recentrings and about \qty{4}{\milli\second} for the day.
\item Each centring clears the whole ladder, and a ladder of millions of ticks no longer fits any cache: the median update
rises from about \qty{29}{\nano\second} to about \qty{135}{\nano\second} at $2^{20}$ and $2^{22}$ ticks.
\item 2~MiB.
\item It is nearly flat, about \qty{30}{\nano\second} from 4\,096 to 65\,536 ticks and 36 at 262\,144: within that range the width costs through
recentrings, not through updates.
\item \textbf{Named result.} On a thirty-percent day of a small-tick stock (300\,000 ticks), a 4\,096-tick ladder recentres
256 times and costs about \qty{13}{\milli\second}, a 65\,536-tick one fifteen times and \qty{4}{\milli\second}, the day's
minimum; a 256-tick ladder recentres almost twenty thousand times and costs \qty{0.7}{\second}.
\item From the instrument's price, tick and the largest daily move expected (a multiple of its volatility): four times that
move in ticks, capped by the memory and cache the book may use.
\item It takes recentring off the critical path: updates continue on the old ladder while the new one is filled, then the
builder switches, at the price of a second ladder and a way to replay the updates made meanwhile.
\item Because a wide ladder rarely recentres or spills into its sparse map, and the tests must exercise those paths.
\item On a large-tick instrument the price travels few ticks and the book is dense near the touch: a small ladder rarely
recentres, and every structure is fast.
\item The count per day and the time each took, with the instrument, so that a width that became too small shows up before it
costs.
\item About \qty{65}{\nano\second} (vector) to \qty{72}{\nano\second} (tree) per event on the small-tick stream whatever the move: no recentring, but twice the ladder's
median every event.
\item Constant-time levels at the price of memory and recentring; it must be protected against the price travelling far,
by its width and by tests on the day that matters.
\end{enumerate}
\end{solution}

\begin{solution}{iq:ll:the-order-book-builder:1}
Look up an order by reference, add to, reduce and remove a level, and read the best level and the first few levels: all in
tens of nanoseconds, every event. Orders in a pool found by \omterm{def:ll:the-order-book-builder:open}{open addressing}; levels in a dense ladder around the touch with a
sparse fallback; invariants checked in debug builds.
\iqlookfor{operations, their frequency and a structure per operation.}
\end{solution}

\begin{solution}{iq:ll:the-order-book-builder:2}
A tree for generality (any price, any tick, no tuning) at logarithmic cost and an allocation per new level; a sorted vector
when the book is shallow and changes concentrate at the touch; an array of levels when the tick is known and the price stays
within a window, with recentring for when it does not. Measure on the instrument's own data.
\iqlookfor{the tick regime and the book's shape decide.}
\end{solution}

\begin{solution}{iq:ll:the-order-book-builder:3}
\omterm{def:ll:the-order-book-builder:open}{Open addressing} with linear probing into flat arrays sized at start-up to twice the live orders, values being slots in a
preallocated pool. Delete by shifting the following entries of the probe run back into the hole when their home slot allows
it, so that no tombstone lengthens later lookups.
\iqlookfor{flat arrays, a pool, and backward-shift deletion.}
\end{solution}

\begin{solution}{iq:ll:the-order-book-builder:4}
A lost or duplicated event (a gap the handler did not flag, an unknown reference), a bug in the builder, or a legitimate cross
(an auction). Mark the book stale, stop trading on it, check the handler's gap counters, and resynchronise from a snapshot;
investigate with the recorded events.
\iqlookfor{stale first, then cause, then recovery.}
\end{solution}

\begin{solution}{iq:ll:the-order-book-builder:5}
Compare it event by event with independent books on shared fixtures, in every language it is written in; exercise its rare
paths (narrow ladders, recentring, snapshots, unknown references); check invariants after every event; replay the busiest and
most volatile days recorded; count allocations.
\iqlookfor{differential testing and the rare paths.}
\end{solution}

\begin{solution}{iq:ll:the-order-book-builder:6}
With daily volatility $\sigma$ in ticks and a recentring at $a$ ticks from the centre, a Brownian price recentres $\sigma^2/a^2$
times a day and a trend of $D$ ticks $D/a$ times. Size the ladder for the largest move expected, recentre in place or in the
background, keep a sparse fallback, and monitor recentrings per instrument.
\iqlookfor{the exit-time argument and a sizing rule.}
\end{solution}
